% ============================================================
% Chapter 7: Results
%   Section: Threats to Validity
% Included from chapters/07-results/results.tex
% ============================================================

\section{Threats to Validity}
\label{sec:threats}

\paragraph{Internal validity.}
\begin{itemize}
    \item \textbf{Labels by a model.} Most labels of the custom set and
    all the candidate labels of the selection experiments were assigned
    by an AI labelling pass. Until the annotator reviews them, a
    result measured against them compares the system with another
    model. Every such result is marked as such.
    \item \textbf{One annotator.} Self-agreement measures whether the
    guide is applied consistently, not whether it is right; a
    systematic misreading applied twice agrees with itself.
    \item \textbf{The selector chooses the set.} The custom set's claims
    are those the system's own selector proposed, so verification is
    measured on the system's choice of claims, not on an independent
    sample of all check-worthy claims.
    \item \textbf{Tuning on the test data.} No threshold has been tuned
    on labelled data yet; the frozen split is there so that none is
    tuned on the claims it is reported on.
\end{itemize}

\paragraph{Construct validity.}
\begin{itemize}
    \item \textbf{Synthetic refuted claims.} The \texttt{FALSE} class of
    the custom set is built by substitution. It tests whether the system
    notices a claim contradicting evidence it can find, not how often or
    how real constructive news is wrong; it is reported apart.
    \item \textbf{Today's web against the claim's date.} The system
    searches without a date limit; the annotator may use only evidence
    available at publication. Newer evidence is flagged, not removed,
    and on X-Fact it dominated (Section~\ref{subsec:pilot-time}).
    \item \textbf{Label mapping.} X-Fact's seven labels are mapped onto
    the system's five, which changes what macro-$F_1$ averages over.
\end{itemize}

\paragraph{External validity.}
\begin{itemize}
    \item \textbf{Size.} With about 100 test claims, a 95\% interval on
    accuracy is roughly $\pm10$ points: only large differences between
    models or conditions will show.
    \item \textbf{Two languages, 37 sources, one period.} The results
    hold for these sources in October 2026, in English and Spanish; the
    selection experiments come from one or two days of feeds.
    \item \textbf{Search availability.} Results depend on which engines
    answered SearXNG on the day; on 13~October Google and Brave refused
    it. Every run records what the health and search checks said
    before it, and claims whose search failed are set apart, but a
    partially answering engine set still shapes the evidence.
    \item \textbf{A moving web.} A search cannot be repeated: the same
    query returns different pages on different days. Strategies are
    compared on the same days, interleaved, rather than across time.
\end{itemize}
